
quotient de \((\underline{O}_{\underline{\underline{G}}_S^{n,r}})^{r+1}\) par le \((r+1)\)-ième facteur. Dans ce système de morphismes, les quotients canoniques se correspondent donc par image inverse. Le système inductif ainsi défini sera noté \(\underline{\underline{G}}_S^{n,\infty}\).

On définit son \(K^{\bullet}\) par la formule :
\begin{equation}\tag{4.9.1}
  K^{\bullet}(\underline{\underline{G}}_S^{n,\infty})
  = \varprojlim_r \bigl(K^{\bullet}(\underline{\underline{G}}_S^{n,r})\bigr) .
\end{equation}

Soient \(\underline{F}_1,\underline{F}_2\) les faisceaux canoniques sur \(\underline{\underline{G}}_S^{n,r}\) (4.5), et \(F_1,F_2\) leur classe dans \(K^{\bullet}(\underline{\underline{G}}_S^{n,r})\). Posons \(c_i=\gamma^i(F_1-n)\). Si on fait varier \(r\), les \(c_i\) se correspondent par les morphismes de transition, et définissent par conséquent des éléments, encore notés \(c_i\), de \(K^{\bullet}(\underline{\underline{G}}_S^{n,\infty})\).

\textbf{Proposition 4.10.} \uline{Soit} \(S\) \uline{un schéma quasi}-\uline{compact}. \uline{Alors}, \uline{avec les notations de} 4.9, \uline{on a un isomorphisme}.~:
\[
  K^{\bullet}(\underline{\underline{G}}_S^{n,\infty})\xrightarrow{\sim} K^{\bullet}(S)[[c_1,\ldots,c_n]] .
\]

D'après 4.8, \(K^{\bullet}(\underline{\underline{G}}_S^{n,r})\) est la \(K^{\bullet}(S)\)-algèbre engendrée par les \(c_i=\gamma^i(F_1-n)\), \(1\le i\le n\), et les \(\gamma^j(F_2-r+n)\), \(1\le j\le r-n\), soumis aux relations déduites de :
\[
  \gamma_t(F_1-n).\gamma_t(F_2-r+n)
  = \gamma_t(r-r)=1,
\]
soit encore :
\[
  \gamma_t(F_2-r+n)=1/\gamma_t(F_1-n).
\]

Cette relation permet d'exprimer les \(\gamma^j(F_2-r+n)\) en fonction des \(c_i\), et donne une série de relations isobares entre les \(c_i\) (\(c_i\) étant supposé de poids \(i\)) en écrivant que les coefficients de \(1/\gamma_t(F_1-n)\) sont nuls en degré \(>r-n\); de plus, ces relations engendrent l'idéal des relations entre les \(c_i\). On peut donc écrire :
